\documentclass[10pt]{article}
\providecommand{\FIGDIR}{fig}
\newcommand{\DIGESTDATE}{22 November 2026}
\input{preamble}

\begin{document}

% ==================================================================== MASTHEAD
\thispagestyle{empty}
\begin{tikzpicture}[remember picture, overlay]
  \fill[Deep] (current page.north west) rectangle
      ($(current page.north east)+(0,-67mm)$);
  \fill[Amber] ($(current page.north west)+(0,-67mm)$) rectangle
      ($(current page.north east)+(0,-68.4mm)$);
  \foreach \x/\y/\r in {22/12/0.9, 41/26/1.5, 63/9/0.7, 88/31/1.1, 112/17/1.9,
                        138/29/0.8, 159/13/1.3, 178/54/1.0, 196/21/1.6,
                        30/44/0.6, 74/57/1.2, 125/48/0.9, 168/6/0.8, 191/59/0.7,
                        52/18/1.0, 100/60/0.6, 148/20/1.1, 66/49/1.4, 205/45/0.8}
    \fill[Sky, opacity=0.30] ($(current page.north west)+(\x mm,-\y mm)$) circle (\r mm);
\end{tikzpicture}

\vspace*{4mm}
{\sffamily\fontsize{7.6}{9}\selectfont\color{Sky}%
 AUTOMATED DAILY LITERATURE DIGEST\;\textbullet\;PhD RESEARCH WATCH\par}
\vspace{2.2mm}
{\sffamily\bfseries\fontsize{27}{29}\selectfont\color{white}%
 Particulate Matter\par}
{\sffamily\fontsize{27}{29}\selectfont\color{Sky}%
 Monitoring \& Health Effects\par}

\vspace{5.0mm}
{\sffamily\fontsize{8.4}{10}\selectfont\color{white}%
 Issue 2026-11-22\;\;\textcolor{Amber}{\textbar}\;\;Entry window 22 November 2026 (full day)\;\;%
 \textcolor{Amber}{\textbar}\;\;10 records screened in scope\par}
\vspace{18mm}

\noindent
\begin{minipage}[t]{0.190\linewidth}\metric{10}{RECORDS IN SCOPE\\(13 EXCLUDED)}{Deep}\end{minipage}\hfill
\begin{minipage}[t]{0.190\linewidth}\metric{6}{OF 10 SUBTOPIC\\BANDS POPULATED}{Teal}\end{minipage}\hfill
\begin{minipage}[t]{0.190\linewidth}\metric{3}{EFFECT ESTIMATES\\WITH A CI}{Sky}\end{minipage}\hfill
\begin{minipage}[t]{0.190\linewidth}\metric{3}{TIER-A\\RECORDS}{Amber}\end{minipage}\hfill
\begin{minipage}[t]{0.190\linewidth}\metric{2}{RECORDS CARRIED ON\\METADATA ALONE}{Clay}\end{minipage}

\vspace{6mm}

\section*{\sffamily\bfseries\color{Deep}Signal of the day}
\vspace{-1.5mm}{\color{Amber}\rule{\linewidth}{1.1pt}}\vspace{2.5mm}

\begin{keybox}
{\sffamily\bfseries\small\color{Deep}In 10,568 women of the China Southwest Birth Cohort, each 10\,$\mu$g\,m$^{-3}$ of whole-pregnancy PM$_{2.5}$ carried 28\,\% higher odds of postpartum depression (OR 1.28, 95\,\% CI 1.06--1.54), rising to 1.79 (1.25--2.57) in women also exposed to tobacco smoke.}\\[1.4mm]
\footnotesize
\textit{Ecotoxicology and Environmental Safety} (Lan, Hu, Xue, Liu). Multicentre prospective cohort; PPD by Edinburgh Postnatal Depression Scale at six weeks (22.1\,\% prevalent), antenatal symptoms by EPDS in the second trimester. Residential PM$_{2.5}$ was assigned for the whole pregnancy and each trimester, and the authors applied two-stage residual inclusion with \emph{wind speed} as an instrument to address endogeneity. The second trimester was the strongest window (OR 1.18, 1.02--1.38) and antenatal depressive symptoms mediated 9.7\,\% of that effect.
\textbf{Where it is weak:} wind speed is not a credible instrument --- it also moves temperature, humidity and every co-pollutant, all plausible direct causes of mood, so the exclusion restriction fails and the IV estimate is not obviously less biased than the naive one; the exposure model itself is never described.
\textbf{Why it matters here:} the tobacco interaction is the interesting claim, and it is exactly the kind of personal-source modification that only personal or home monitoring --- not a residential grid --- can separate from ambient PM$_{2.5}$.\par
\end{keybox}

\vspace{3mm}

\section*{\sffamily\bfseries\color{Deep}What changed today}
\vspace{-1.5mm}{\color{Amber}\rule{\linewidth}{1.1pt}}\vspace{2.5mm}

\footnotesize
\begin{itemize}

\item \textbf{PM$_1$ looks worse than PM$_{2.5}$, on evidence that cannot carry the comparison.} An umbrella review of observational meta-analyses (\textit{Ecotox.\ Environ.\ Saf.}, GRADE-graded) reports PM$_1$ relative risks of 1.30 for metabolic syndrome, 1.41 childhood obesity, 1.61 childhood hypertension and 1.59 NAFLD, above the PM$_{2.5}$ and PM$_{10}$ figures. No CIs and no increments are given, and PM$_1$ in these cohorts is modelled and near-collinear with PM$_{2.5}$ --- so the ``size effect'' may be increment scaling, not biology. It still matters operationally: optical low-cost sensors report a PM$_1$ channel, and this is the literature that will be cited to justify using it.

\item \textbf{Brake wear: compliant on the regulatory cycle, non-compliant off it.} Dynamometer tests of two pad materials (\textit{J.\ Hazard.\ Mater.}) stayed inside the Euro 7 PM$_{10}$ limit under the WLTP-Brake cycle but exceeded it under other driving conditions, and particle number tracked the history of preceding brake events. A Transformer on time-series inputs predicted particle number at R$^2$ 0.85\,$\pm$\,0.03, with SHAP putting brake temperature and brake power first. The cycle, not the pad, is doing the regulatory work.

\item \textbf{Satellite gap-filling in Mexico, at 3.7--9.9\,$\mu$g\,m$^{-3}$ RMSE.} MAIAC AOD plus ERA5 reanalysis, SINAICA monitors and OpenStreetMap covariates in mixed-effects and GAM models give daily PM$_{2.5}$ for 2004--2022 in four metropolitan areas (\textit{PeerJ}); Mexico City's yearly leave-one-out RMSE runs 3.67--9.91. Validation is leave-one-out over a sparse monitor set --- the same optimism the 21 November issue flagged at 1.45$\times$ --- and PM$_{10}$-derived proxy PM$_{2.5}$ carries an unquantified ratio assumption.

\item \textbf{Two experimental exposures at opposite ends of plausibility.} A wildland-firefighter model (bioRxiv preprint) ran Apoe$^{-/-}$ mice at 40\,mg\,m$^{-3}$ Douglas-fir smoke for 8--16 weeks, dose-matched by cumulative deposition to 7--14 years of service: eNOS uncoupling with endothelial dysfunction maximal by 8 weeks, while aortic stiffening and blood pressure kept rising to 16 weeks. Separately, gestational PM$_{2.5}$ at only 10\,$\mu$g\,m$^{-3}$ with 39\,$^\circ$C heat, 2\,h/day, reportedly thickened the ventricular wall and impaired systolic function in day-old mice via myo-inositol and \textit{Ucp2} loss --- a large structural effect for a near-ambient dose, and the arm that most needs litter-level reanalysis.

\item \textbf{Exposure basis (x1).} Today's three health-side estimates all rest on \emph{modelled} residential PM$_{2.5}$; the measured exposures in the issue are both laboratory (smoke chamber, brake dynamometer), and the only network-grade measurement is the Mexican monitor-plus-satellite fusion. No low-cost-sensor record today --- the first such gap since 21 November, when the single sensor record was the issue's lead.

\end{itemize}

\vspace{2mm}

\begin{keybox}
{\sffamily\bfseries\footnotesize\color{Deep}Window continuity}\\[1.2mm]
{\fontsize{7.2}{8.8}\selectfont
The 21 November issue closed the day before. This issue was built on the \textbf{23 November 13:25 UTC (07:25 CST) run} (the 05:05 UTC slot fired about eight hours late). Before 22:00 CST the newest buildable day is yesterday, so this run owes \textbf{22 November} alone --- a Sunday, no rollup due; \textbf{22 November is covered in full}. W48 (22--28 November) falls due on Saturday 28 November and the November monthly on 30 November. Seven PubMed connector PMIDs with entry date 21 November, deposited after that issue's harvest, were folded into today's candidate set and dropped by the \texttt{seen.json} and corpus check or screened here. No \textsc{backfill} records today. OpenAlex returned HTTP 403 again; the alphaXiv, bioRxiv/medRxiv and Wiley Scholar Gateway connectors were not available on this session, and Scite and Consensus were not run.\par}
\end{keybox}

\vspace{3mm}
\par\ifdim\dimexpr\pagegoal-\pagetotal\relax<18\baselineskip\newpage\fi
\section*{\sffamily\bfseries\color{Deep}Corpus analytics}
\vspace{-1.5mm}{\color{Amber}\rule{\linewidth}{1.1pt}}\vspace{2.5mm}

\noindent\begin{minipage}{\linewidth}\includegraphics[width=\linewidth]{\FIGDIR/f1_subtopics.png}
\figcap{Subtopic assignment of the 10 records across six of ten bands: reproductive \& developmental 3, occupational \& indoor 2, cardiovascular \& metabolic 2, and one each for neuro/mental health, exposure assessment and sensing. Respiratory, burden, mechanistic and other clinical are empty today.}\end{minipage}

\vspace{2mm}
\noindent\begin{minipage}{\linewidth}\includegraphics[width=\linewidth]{\FIGDIR/f2_design_endpoint.png}
\figcap{Left --- architecture: two prospective/retrospective cohorts, two animal exposure studies, one source-isolated emission test, one satellite--surface fusion, two reviews and two metadata-only records. Right --- endpoints are singletons apart from the metadata-only bar.}\end{minipage}

\vspace{2mm}
\noindent\begin{minipage}{\linewidth}\includegraphics[width=\linewidth]{\FIGDIR/f3_metric_geography.png}
\figcap{Left --- particle metric as classified from the record fields: PM$_{2.5}$ dominates, with one PM$_1$/PM$_{2.5}$/PM$_{10}$ comparison and one brake-wear PM$_{10}$/particle-number record. Right --- geography by record: China 4 (one cohort, one metadata-only record and two laboratory studies), global / multi-region 3 (the PM$_1$ umbrella review, the multi-country Asian household-air-pollution review and one record with no stated site), North America 2 (the four-state US APWA cohort and the US firefighter smoke model) and Latin America 1 (Mexico).}\end{minipage}

\vspace{2mm}
\noindent\begin{minipage}{\linewidth}\includegraphics[width=\linewidth]{\FIGDIR/f6_lifecourse.png}
\figcap{Life-course windows. Preconception / in utero carries four records (postpartum depression, the APWA abruption cohort, the gestational PM$_{2.5}$--heat mouse study and the household-air-pollution review of reproductive outcomes); infancy and early childhood one (the PM$_1$ umbrella review's obesity, hypertension, asthma and autism outcomes); working-age adulthood two (the firefighter smoke model and battery-materials workers). No adolescent or older-adult record today.}\end{minipage}

\vspace{2mm}
\noindent\begin{minipage}{\linewidth}\includegraphics[width=\linewidth]{\FIGDIR/f4_heatmap.png}
\figcap{Subtopic $\times$ study architecture. Reproductive \& developmental spans a cohort, an animal study and a review --- the widest band today; occupational \& indoor pairs a review with a metadata-only occupational record; sensing holds the brake-dynamometer test and exposure assessment the satellite fusion.}\end{minipage}

\vspace{2mm}

\IfFileExists{\FIGDIR/extra_figs.tex}{\input{\FIGDIR/extra_figs.tex}}{}



\noindent\begin{minipage}{\linewidth}\includegraphics[width=\linewidth]{\FIGDIR/f5_forest.png}
\figcap{Ratio scale, three estimates from one study. All come from the China Southwest Birth Cohort (n\,=\,10,568) and one outcome, postpartum depression: whole-pregnancy PM$_{2.5}$ OR 1.28 (1.06--1.54) per 10\,$\mu$g\,m$^{-3}$, second trimester 1.18 (1.02--1.38), and 1.79 (1.25--2.57) in the tobacco-smoke-exposed subgroup. This is a within-study pattern, not a between-study comparison: the subgroup estimate shares its data with the overall one, so no heterogeneity statistic is meaningful and nothing is pooled.}\end{minipage}

\vspace{4mm}

\par\vskip 0pt plus 24\baselineskip\penalty-200\vskip 0pt plus -24\baselineskip
\section*{\sffamily\bfseries\color{Deep}Paper-level digest}
\vspace{-1.5mm}{\color{Amber}\rule{\linewidth}{1.1pt}}\vspace{2.5mm}

{\fontsize{7.4}{9}\selectfont\color{Slate}Ordered by cluster. Each entry gives the
methodological core and the finding that matters, not an abstract paraphrase. Records
marked \textit{metadata only} carry no recoverable abstract and assert no finding.\par}

\band{Neuro / mental health\hfill 1 record}{Deep}

\paperentry{China Southwest Birth Cohort, n=10,568: whole-pregnancy PM$_{2.5}$ OR 1.28 (1.06-1.54) per 10 $\mu$g\,m$^{-3}$ for postpartum depression; 1.79 among tobacco-exposed; 9.7\% mediated by antenatal symptoms}
{Prenatal ambient PM(2.5) exposure during pregnancy and postpartum depression: A prospective Birth Cohort study in China.}
{Ecotoxicol Environ Saf}
{Multicentre China Southwest Birth Cohort, 10,568 women; PPD by EPDS at 6 weeks postpartum (22.1\% prevalent), antenatal symptoms by EPDS in the second trimester. Each 10 $\mu$g\,m$^{-3}$ of whole-pregnancy PM$_{2.5}$ was associated with OR 1.28 (95\% CI 1.06-1.54) for PPD after a two-stage residual-inclusion IV correction using wind speed; OR 1.79 (1.25-2.57) among tobacco-smoke-exposed women (P-heterogeneity 0.046); second trimester OR 1.18 (1.02-1.38), increment not restated. Antenatal symptoms mediated 9.7\% of the second-trimester effect. Weakness: wind speed is a weak and probably invalid instrument - it moves heat, humidity and other pollutants that plausibly affect mood - and the exposure model is not described. Relevance: modest, but a large cohort where personal or home sensors would test whether the IV correction removes or adds bias.}
{10.1016/j.ecoenv.2026.121090}

\par\vskip 0pt plus 10\baselineskip\penalty-200\vskip 0pt plus -10\baselineskip
\band{Cardiovascular \& metabolic\hfill 2 records}{Teal}

\paperentry{Umbrella review of PM$_{1}$/PM$_{2.5}$/PM$_{10}$ meta-analyses with GRADE: PM$_{1}$ RR 1.30 metabolic syndrome, 1.41 childhood obesity, 1.61 childhood hypertension, 1.59 NAFLD - higher than PM$_{2.5}$/PM$_{10}$}
{Different diameters of particulate matter exposure and health outcomes: An umbrella review.}
{Ecotoxicol Environ Saf}
{Umbrella review of observational meta-analyses, harmonising effect estimates by particle size and grading certainty with GRADE. PM$_{1}$ carried the largest pooled relative risks: metabolic syndrome RR 1.30, childhood obesity 1.41, childhood hypertension 1.61 and NAFLD 1.59 (moderate certainty); similar size gradients were reported for autism and asthma. CIs, increments and the number of included meta-analyses are not given in the abstract, so these are not plotted. Weakness: PM$_{1}$ meta-analyses draw on few, mostly Chinese cohorts with modelled PM$_{1}$ highly collinear with PM$_{2.5}$, and RRs per unspecified increments cannot be compared across fractions - the "size effect" may be an artefact of increment scaling. Relevance: direct - many optical low-cost sensors report PM$_{1}$, but its health signal is only as good as its separation from PM$_{2.5}$.}
{10.1016/j.ecoenv.2026.120993}

\paperentry{Apoe-/- mice, Douglas-fir smoke 40 mg/m3, 2 h/d, 5 d/wk, 8 or 16 wk (dose matched to 7-14 y of wildland firefighting): eNOS uncoupling, endothelial dysfunction maximal by 8 wk, aortic stiffening still worsening at 16 wk}
{Occupationally Relevant Wildfire Smoke Inhalation Impairs Nitric Oxide Signaling and Promotes Progressive Aortic Stiffening in Hypercholesterolemic Mice.}
{bioRxiv}
{Preprint, not peer reviewed. Male Apoe-/- mice inhaled Douglas fir smoke or filtered air at a 40 mg/m3 PM target for 8 or 16 weeks. Smoke raised circulating TNF-alpha, IL-1beta and IL-6 with diffuse aortic NF-kB activation, disrupted endothelial adherens junctions, raised ICAM-1, VCAM-1 and MCP-1 with monocyte recruitment, and increased NOX2, superoxide and eNOS uncoupling with lipid peroxidation and intimal apoptosis. A pro-fibrotic shift (TGF-beta1, myofibroblasts, medial and adventitial collagen) accompanied falling cyclic distensibility, attenuated NO-mediated vasodilation and higher systolic and diastolic pressure. Endothelial dysfunction plateaued by 8 weeks while stiffening progressed to 16 weeks. No effect sizes or CIs are reported in the abstract. Weakness: 40 mg/m3 is two to three orders above field firefighter PM$_{2.5}$ and the dose match is by cumulative deposition, not concentration, so the rate-dependent biology may not transfer; Apoe-/- mice are already hyperlipidaemic. Relevance: supplies the mechanism that personal-monitor firefighter studies cannot, and sets the cumulative-dose metric those monitors would have to integrate.}
{10.64898/2026.05.18.725908}

\par\vskip 0pt plus 10\baselineskip\penalty-200\vskip 0pt plus -10\baselineskip
\band{Reproductive \& developmental\hfill 3 records}{Sage}

\paperentry{APWA: 13.6 million linked deliveries in four US states (2000-2016), 144,930 abruptions; daily modelled PM$_{2.5}$, NO$_2$, O$_3$ and weather from 90 days pre-LMP to 12 months postpartum}
{Ambient Air Pollution, Weather and Placental Abruption (APWA): Design, Data Linkage and Preliminary Findings.}
{Paediatr Perinat Epidemiol}
{Cohort-profile paper. Vital records for live births and fetal deaths in California, Florida, Massachusetts and South Carolina (2000-2016) linked to delivery hospitalisations: 13,612,485 of 14,756,201 births (94\%) linked to abruption diagnosis and to daily modelled pollutant and weather exposures. Abruption prevalence 1.0-1.4\% by state (86,803 in California). Exposure distributions were similar for abruption and non-abruption births, with heterogeneity by state and season. No effect estimate yet. Weakness: residence at delivery only, so pre-conception and early-pregnancy windows - the paper's own chronic-abruption hypothesis - carry the largest misclassification; exposure model not named in the abstract. Relevance: the acute vs chronic window contrast will hinge on daily exposure accuracy, which is where dense networks add most over gridded models.}
{10.1111/ppe.70211}

\paperentry{Pregnant mice, PM$_{2.5}$ 10 $\mu$g\,m$^{-3}$ and/or 39 C, 2 h/d E6.5-E18.5: co-exposure thickened ventricular wall and impaired systolic function at P1, with lower myo-inositol, Ucp2 and ATP}
{Gestational co-exposure to PM(2.5) and heat impaired cardiac development in offspring.}
{Ecotoxicol Environ Saf}
{Pregnant mice exposed to filtered air, PM$_{2.5}$ (10 $\mu$g\,m$^{-3}$), heat (39 C) or both. Co-exposure increased ventricular wall thickness and impaired systolic function on postnatal day 1, with inflammatory infiltration, oxidative damage and mitochondrial ultrastructural change. Paired transcriptomics and metabolomics implicated inositol phosphate and pantothenate/CoA metabolism; myo-inositol and Ucp2 fell, as did ATP, while DRP1 and FIS1 protein rose and cardiac development genes were downregulated. No effect sizes, CIs or litter-level n are given in the abstract. Weakness: a 2 h/day exposure at an ambient-like 10 $\mu$g\,m$^{-3}$ producing structural cardiac change is a surprisingly large effect for the dose and invites a litter-clustering and multiplicity check; the 39 C arm is severe heat stress, so the "co-exposure" interaction may be heat-driven. Relevance: indirect - but the PM$_{2.5}$-heat interaction is the case for co-locating temperature with PM on every sensor node.}
{10.1016/j.ecoenv.2026.121081}

\paperentry{Authors reply on interpreting differential susceptibility to prenatal PM$_{2.5}$ exposure}
{Reply to 'Interpreting Differential Susceptibility to Prenatal PM(2.5) Exposure'.}
{Respirology}
{\textsc{metadata only.} Metadata only: no abstract was retrievable from the harvest legs, Crossref or Semantic Scholar. Authors reply (Respirology; Huang, Pan, Yao) to a letter on interpreting differential susceptibility to prenatal PM$_{2.5}$ exposure. No finding is asserted here, and the original study is not identified in the record. Relevance: low - carried because effect-modification disputes in prenatal PM$_{2.5}$ work usually turn on exposure-model resolution, which is the question dense networks address.}
{10.1002/resp.70342}

\par\vskip 0pt plus 10\baselineskip\penalty-200\vskip 0pt plus -10\baselineskip
\band{Exposure assessment \& modelling\hfill 1 record}{Clay}

\paperentry{Four Mexican metropolitan areas, 2004-2022: MAIAC AOD + ERA5 + OSM covariates in mixed-effects and GAM models give daily PM$_{2.5}$ with LOOCV RMSE 3.67-9.91 $\mu$g\,m$^{-3}$ in Mexico City}
{PM(2.5) modeling in Mexican cities using Aerosol optical depth and Meteorological reanalysis data.}
{PeerJ}
{Hybrid framework combining linear mixed-effects and generalised additive models to produce spatially continuous daily PM$_{2.5}$ for 2004-2022 in four Mexican metropolitan areas, including days and places without satellite retrievals. In the best region, the Mexico City Metropolitan Area, yearly leave-one-out RMSE ranged 3.67-9.91 $\mu$g\,m$^{-3}$; station-based and reanalysis meteorological inputs were compared, and reanalysis inputs sustained estimates in cities usually excluded for sparse monitoring. Weakness: leave-one-out over a sparse monitor set is close to the random-CV inflation flagged on 21 November - no spatial-block or leave-one-station-out figure is reported - and PM$_{10}$-derived proxy PM$_{2.5}$ propagates an unquantified ratio assumption. Relevance: direct - this is the low-resource-setting gap-filling problem a low-cost network is meant to close, and the RMSE here is the bar such a network has to beat.}
{10.7717/peerj.21845}

\par\vskip 0pt plus 10\baselineskip\penalty-200\vskip 0pt plus -10\baselineskip
\band{Sensing, forecasting \& instrumentation\hfill 1 record}{Amber}

\paperentry{Brake dynamometer, two pad materials: PM$_{10}$ within the Euro 7 WLTP-Brake limit but exceeded under other cycles; Transformer predicts particle number from brake time series at R2 0.85 +/- 0.03}
{Modeling brake wear particle emissions under diverse driving conditions: A data-driven time-series approach.}
{J Hazard Mater}
{Two pad materials from different vehicles tested on a brake dynamometer across multiple driving conditions. PM$_{10}$ complied with the forthcoming Euro 7 limit under the WLTP-Brake cycle but exceeded it under other cycles, and particle-number emissions varied between cycles in a way that tracked preceding brake events, showing strong temporal dependence. With time-series inputs a Transformer predicted particle number best (R2 0.85 +/- 0.03), XGBoost acceptably; SHAP attributed high-emission tests to brake temperature and brake power, with speed positively related and amplified by power. Weakness: two pad formulations on one dynamometer, no on-road validation, and R2 on held-out cycles from the same rig overstates transferability; no PM$_{1}$ or composition. Relevance: direct - non-exhaust PM is the growing urban source, and a temperature/power-driven emission model is what a roadside low-cost network would need to attribute its own spikes.}
{10.1016/j.jhazmat.2026.144288}

\par\vskip 0pt plus 10\baselineskip\penalty-200\vskip 0pt plus -10\baselineskip
\band{Occupational \& indoor\hfill 2 records}{Amber}

\paperentry{Systematic review, 28 studies: household solid-fuel smoke in Asian women linked to respiratory symptoms (aOR 3.0), birth defects (aOR up to 7.09), stillbirth (aOR 1.34)}
{Indoor air pollution and health outcomes in women in developing Asian countries: a systematic review.}
{Environ Anal Health Toxicol}
{PRISMA review of five databases, 28 studies of women in developing Asian countries, quality appraised with NOS, AXIS and OHAT. Household air pollution, mainly solid-fuel combustion, was associated with respiratory symptoms (aOR 3.0), reduced lung function, birth defects (aOR up to 7.09), stillbirth (aOR 1.34), cardiometabolic risk, genotoxicity and cognitive decline. No pooling and no CIs in the abstract. Weakness: exposure is overwhelmingly fuel type rather than measured PM, the selected maxima (7.09) are cherry-picked single studies, and confounding by poverty is the obvious alternative. Relevance: household-energy studies are where cheap personal PM monitors would replace the fuel-type proxy.}
{10.5620/eaht.2026012}

\paperentry{Occupational lithium NMC oxide exposure and lower ventilation function, with inflammatory and transcriptomic mediation}
{Association between lithium nickel manganese cobalt oxide exposure and lower pulmonary ventilation function indices: the indirect effects of blood cell-derived systemic inflammatory markers and transcriptomic profiles}
{Environment International}
{\textsc{metadata only.} Metadata only: no abstract was retrievable from the harvest legs, Crossref or Semantic Scholar. The record (Environment International; Li, Liu, Hu, Fang) relates lithium nickel manganese cobalt oxide exposure to lower pulmonary ventilation function and examines mediation by blood-cell-derived inflammatory markers and transcriptomic profiles. No finding, exposure metric, cohort or effect size is asserted here. Relevance: moderate - battery-materials workplaces are an emerging source of metal-rich respirable particles that mass-only optical sensors cannot distinguish from background aerosol.}
{10.1016/j.envint.2026.110608}

\clearpage
\section*{\sffamily\bfseries\color{Deep}Corpus register}
\vspace{-1.5mm}{\color{Amber}\rule{\linewidth}{1.1pt}}\vspace{2.5mm}

\input{table_rows}

\vspace{2mm}

\begin{keybox}
{\sffamily\bfseries\footnotesize\color{Deep}Provenance and method}\\[1.2mm]
{\fontsize{7.2}{8.8}\selectfont
Entry window \textbf{22 November 2026}, single day (Sunday). Harvester legs: \textbf{PubMed} health \textbf{8}, sensing \textbf{2}; \textbf{Europe PMC} \textbf{5} (3 AGRICOLA retro-index records dropped); \textbf{Crossref by ISSN} \textbf{7}; \textbf{arXiv} no in-window entry; \textbf{OpenAlex} HTTP 403 on the harvester leg. 19 raw, 17 unique after cross-source dedupe, 0 already carried. PubMed connector (EDAT 22 November, health and instrumentation axes) 15 PMIDs (\texttt{\_conn\_1123.py}), including 7 with entry date 21 November deposited after that harvest: 1 already seen, 1 merged onto a cached record, 6 appended --- \textbf{23 screened}. Abstracts: efetch and Crossref filled 14 of 15 attempts; 2 records carried on metadata alone (one Elsevier deposit, one correspondence).
alphaXiv, bioRxiv/medRxiv and Wiley Scholar Gateway connectors were not available on this session; \textbf{Scite} retraction check and Consensus not run. No Amass or Consensus literature backfill this run.
\textbf{10 admitted}, \textbf{13 rejected} with logged reasons (\texttt{\_rej\_1123\_22.py}). Affiliations not retrieved this issue. Trial watch 22--23 November: ClinicalTrials.gov 0 hits in window; Amass TrialCore returned 25 records by last-update date, every on-topic one already carried (REACT, AIRWISE, the Leicester hospital-environment study, the AQAS parent survey) --- \textbf{0 new}; \texttt{trials.json} now 31 trials across 49 windows.
Abstracts remain the copyright of their publishers.\par}
\end{keybox}

\end{document}
